\section{Induction Heads: el circuito que implementa el In-Context Nearest Neighbor}

Volvamos ahora al phase-change mystery. Una induction head procesa el patrón
\[
[A][B]\ \cdots\ [A]\quad\longrightarrow\quad[B].
\]
Cuando el token actual es $A$, busca una $A$ que haya aparecido antes, lee la $B$ que sigue a esa $A$ anterior y eleva la next-token probability de $B$.

\subsection{Induction pattern}

El mecanismo suele requerir una composition en dos pasos. La previous-token head escribe en la posición $j$ la token identity de la posición $j-1$; en la capa siguiente, la induction head hace un QK match entre la $A$ de la posición actual y la previous-token feature que lleva cada posición candidata, y mediante un copying OV emite el token actual $B$ de la posición candidata.

\lecturefigure{slide-49-induction-pattern.jpg}{Induction pattern: se busca un prefix idéntico anterior y se copia el token que lo sigue.}{Video oficial de Stanford Online 00:46:08--00:47:01.}

\begin{lstlisting}[caption={Pseudocódigo conceptual del induction-head algorithm.}]
def induction_prediction(tokens, current_index):
    current_token = tokens[current_index]
    candidates = []
    for source in range(1, current_index):
        if tokens[source - 1] == current_token:
            candidates.append(tokens[source])
    return learned_weighted_copy(candidates)
\end{lstlisting}

\begin{importantbox}{In-context nearest neighbor}
La induction head trata el context como una base de datos temporal: la query es el pattern actual, las keys son los patterns históricos y el value es el token que apareció después de cada pattern histórico. Se parece a k-NN, pero la similitud, la representation, las rutas de composición y la output transformation las aprende la red.
\end{importantbox}

\teachervoice{La definición condensada de Olah es: buscar la previous copy, mirar una posición hacia adelante y predecir que lo que ocurrió la vez anterior volverá a ocurrir. Las heads reales pueden emparejar frases, token classes o vecinos semánticos, no solo una $A$ literalmente idéntica.}

\subsection{Por qué es el workhorse del meta-learning}

Una copying head simple solo eleva a ciegas los tokens repetidos en muchas posiciones posibles; el induction circuit, en cambio, primero identifica «qué situación pasada se parece a la actual» y después copia la continuation de esa situación. Este conditional retrieval se parece más a un context algorithm generalizable, por lo que sirve para texto repetido, imitación de formato, aprendizaje local de funciones y correspondencias de traducción.

\lecturefigure{slide-50-induction-workhorse.jpg}{Induction heads: el workhorse del meta-learning, que combina pattern matching y copying.}{Video oficial de Stanford Online 00:47:02--00:47:05.}

\lecturefigure{slide-51-multiple-induction-heads.jpg}{En el modelo existen varias induction heads, que cubren distintos patterns y logit effects.}{Video oficial de Stanford Online 00:47:06--00:47:29.}

\subsection{Del attention pattern a la matrix signature}

En una induction head clásica, el OV debe tener una copying tendency y el QK debe preferir el same-token/same-pattern match sobre la shifted feature de la primera capa. Por ello pueden usarse QK/OV traces o eigenvalue summaries para ubicar las heads en un diagrama bidimensional: las heads de la esquina superior derecha, que combinan strong QK matching y OV copying, son induction candidates.

\lecturefigure{slide-52-head-taxonomy.jpg}{Clasificación espectral y por patrones de induction heads, otras heads e induction-ish heads.}{Video oficial de Stanford Online 00:47:30--00:47:47.}

\lecturefigure{slide-53-qk-ov-trace-scatter.jpg}{QK/OV trace scatter: dos circuit signatures para localizar induction candidates.}{Video oficial de Stanford Online 00:47:48--00:48:15.}

\begin{warningbox}{La signature no es la definición}
Que una head caiga en cierta región del scatter solo indica que sus matrices tienen propiedades de candidata. Para establecer un verdadero induction mechanism se necesitan además el attention pattern, la path composition, ejemplos y ablation; los modelos modernos pueden lograr un comportamiento similar mediante otras bases, distributed features o la cooperación de varias heads.
\end{warningbox}

\lecturefigure{slide-54-induction-meta-learning-summary.jpg}{Principal aportación de las induction heads al meta-learning e idea para detectarlas automáticamente.}{Video oficial de Stanford Online 00:48:16--00:48:21.}

\lecturefigure{slide-55-other-local-heads.jpg}{Otras heads suelen ser más local y realizan token lookup o position tracking.}{Video oficial de Stanford Online 00:48:22--00:48:33.}

\lecturefigure{slide-56-lookup-similar-context.jpg}{El núcleo de la induction head: consultar qué ocurrió después en una situación histórica similar.}{Video oficial de Stanford Online 00:48:34--00:48:41.}

\subsection{De vuelta a la phase-change hypothesis}

La viabilidad estructural ya está establecida: una composition de dos capas puede implementar un conditional lookup, y las induction candidates pueden identificarse a partir del QK/OV y del attention pattern. Con ello, la hipótesis de investigación se vuelve concreta: la ICL phase change es el momento en que el modelo descubre durante el entrenamiento los induction-head circuits.

\lecturefigure{slide-57-phase-change-hypothesis.jpg}{Hipótesis central: la formación de induction heads impulsa la ICL phase change.}{Video oficial de Stanford Online 00:48:42--00:49:53.}

\teachervoice{El ponente vuelve a usar la palabra hypothesis y enseguida dice «let's check it». Esto refleja el ritmo correcto de la explicación mecanicista: primero se plantea, a partir de la estructura, una hipótesis falsable y después se buscan intervenciones y evidencia en distintos modelos.}

\subsection{Resumen de la sección}

El induction circuit combina la previous-token information con una copying head de la capa siguiente para formar el algoritmo $[A][B]\ldots[A]\to[B]$. Es un context nearest-neighbor mechanism aprendible, más preciso que el copying incondicional, y ofrece una explicación candidata estructuralmente creíble de la phase change.
